\begin{figure}[H]
    \centering
    \begin{tikzpicture}[x=1.0cm, y=1.0cm]

        \newcommand{\distframe}[3]{%
            \begin{scope}[shift={(#1,#2)}]
                \draw[draw=figmute, line width=0.4pt] (0,0) rectangle (3.05,1.30);
                \node[panel, font=\scriptsize\bfseries, anchor=south west] at (0,1.36) {#3};
            \end{scope}}

        % uniform
        \distframe{0}{0}{\tl{uniform}}
        \begin{scope}[shift={(0,0)}]
            \foreach \i in {0,...,60} {
                \pgfmathsetmacro{\xx}{0.05+\i*0.049}
                \pgfmathsetmacro{\yy}{0.10+1.10*rnd}
                \fill[figaieD, opacity=0.75] (\xx,\yy) circle (0.60pt);
            }
        \end{scope}

        % trimodal
        \distframe{3.55}{0}{\tl{trimodal}}
        \begin{scope}[shift={(3.55,0)}]
            \foreach \i in {0,...,60} {
                \pgfmathsetmacro{\xx}{0.05+\i*0.049}
                \pgfmathsetmacro{\pick}{int(random(0,2))}
                \pgfmathsetmacro{\ctr}{ifthenelse(\pick==0,0.20,ifthenelse(\pick==1,0.50,0.80))}
                \pgfmathsetmacro{\yy}{0.10+1.10*(\ctr+0.035*(rnd-0.5)*2)}
                \fill[figaieD, opacity=0.75] (\xx,\yy) circle (0.60pt);
            }
        \end{scope}

        % sorted
        \distframe{5.325}{-2.15}{\tl{sorted}}
        \begin{scope}[shift={(5.325,-2.15)}]
            \draw[figaieD, line width=0.9pt] (0.05,0.12) -- (3.00,1.20);
            \node[note, font=\tiny, anchor=north west, text=figdrmD] at (0.06,-0.06)
                {δυσμενέστερη για σωρό};
        \end{scope}

        % reverse
        \distframe{1.775}{-2.15}{\tl{reverse}}
        \begin{scope}[shift={(1.775,-2.15)}]
            \draw[figaieD, line width=0.9pt] (0.05,1.20) -- (3.00,0.12);
            \node[note, font=\tiny, anchor=north west, text=figregD] at (0.06,-0.06)
                {ευνοϊκότερη για σωρό};
        \end{scope}

        % adversarial
        \distframe{7.10}{0}{\tl{adversarial}}
        \begin{scope}[shift={(7.10,0)}]
            \foreach \i in {0,...,60} {
                \pgfmathsetmacro{\xx}{0.05+\i*0.049}
                \pgfmathsetmacro{\yy}{0.10+0.55*rnd}
                \fill[figaieD, opacity=0.55] (\xx,\yy) circle (0.60pt);
            }
            \draw[dashed, draw=figmute, line width=0.35pt] (0.05,0.68) -- (3.00,0.68);
            \foreach \s in {1,...,7} {
                \pgfmathsetmacro{\xx}{0.05+\s*0.40}
                \pgfmathsetmacro{\yy}{0.70+0.52*\s/7}
                \draw[figdrmD, line width=0.35pt, opacity=0.55] (\xx,0.12) -- (\xx,\yy);
                \fill[figdrmD] (\xx,\yy) circle (0.8pt);
            }
        \end{scope}


    \end{tikzpicture}
    \caption{Οι πέντε κατανομές εισόδου, με την τιμή κάθε στοιχείου συναρτήσει της
    θέσης του στον πίνακα. Η \tl{uniform} και η \tl{trimodal} δεν έχουν καμία
    διάταξη: διαφέρουν μόνο ως προς το ποιες τιμές εμφανίζονται, με τη δεύτερη να
    συγκεντρώνει τις τιμές γύρω από τρία κέντρα. Οι δύο διατεταγμένες κατανομές
    είναι τα άκρα για την επιλογή με σωρό, καθώς στην αύξουσα κάθε νέο στοιχείο
    υπερβαίνει τη ρίζα ενώ στη φθίνουσα κανένα δεν την υπερβαίνει μετά τα πρώτα
    \tl{k}. Η \tl{adversarial} διατηρεί τη δυσμενή συμπεριφορά χωρίς να είναι
    ταξινομημένη: το σώμα των τιμών παραμένει στο κάτω ήμισυ του εύρους, ενώ σε
    τακτά διαστήματα παρεμβάλλονται μεμονωμένες αιχμές των οποίων η τιμή αυξάνει
    γραμμικά κατά μήκος του πίνακα· η διακεκομμένη γραμμή σημειώνει τη μέση τιμή
    του εύρους.}
    \label{fig:meth-dists}
\end{figure}
